\documentclass[aps,preprint]{revtex4}%
\usepackage{amsfonts}%
\usepackage{amsmath}%
\setcounter{MaxMatrixCols}{30}%
\usepackage{amssymb}%
\usepackage{graphicx}
%TCIDATA{OutputFilter=latex2.dll}
%TCIDATA{Version=4.10.0.2363}
%TCIDATA{CSTFile=revtex4.cst}
%TCIDATA{Created=Monday, June 28, 2004 22:18:25}
%TCIDATA{LastRevised=Monday, June 28, 2004 22:19:16}
%TCIDATA{<META NAME="GraphicsSave" CONTENT="32">}
%TCIDATA{<META NAME="DocumentShell" CONTENT="Articles\SW\REVTeX 4">}
\newtheorem{theorem}{Theorem}
\newtheorem{acknowledgement}[theorem]{Acknowledgement}
\newtheorem{algorithm}[theorem]{Algorithm}
\newtheorem{axiom}[theorem]{Axiom}
\newtheorem{claim}[theorem]{Claim}
\newtheorem{conclusion}[theorem]{Conclusion}
\newtheorem{condition}[theorem]{Condition}
\newtheorem{conjecture}[theorem]{Conjecture}
\newtheorem{corollary}[theorem]{Corollary}
\newtheorem{criterion}[theorem]{Criterion}
\newtheorem{definition}[theorem]{Definition}
\newtheorem{example}[theorem]{Example}
\newtheorem{exercise}[theorem]{Exercise}
\newtheorem{lemma}[theorem]{Lemma}
\newtheorem{notation}[theorem]{Notation}
\newtheorem{problem}[theorem]{Problem}
\newtheorem{proposition}[theorem]{Proposition}
\newtheorem{remark}[theorem]{Remark}
\newtheorem{solution}[theorem]{Solution}
\newtheorem{summary}[theorem]{Summary}
\newenvironment{proof}[1][Proof]{\noindent\textbf{#1.} }{\ \rule{0.5em}{0.5em}}
\begin{document}
\preprint{HEP/123-qed}
\title[Short title for running header]{Long title that appears on the title page}
\author{First Author}
\affiliation{My Institution}
\author{Second Author}
\affiliation{My Institution}
\author{Third Author}
\affiliation{Other Institution}
\keywords{one two three}
\pacs{PACS number}

\begin{abstract}
Shell document for REV\TeX{} 4.

\end{abstract}
\volumeyear{year}
\volumenumber{number}
\issuenumber{number}
\eid{identifier}
\date[Date text]{date}
\received[Received text]{date}

\revised[Revised text]{date}

\accepted[Accepted text]{date}

\published[Published text]{date}

\startpage{101}
\endpage{102}
\maketitle
\tableofcontents


For convenience we can express $\mathbf{\theta}$ as a diagonal matrix
\begin{equation}
\mathbf{\hat{\theta}=}\left(
\begin{array}
[c]{cccc}%
0 & 0 & \cdots & 0\\
\vdots & \theta_{2} &  & \vdots\\
\vdots &  & \ddots & 0\\
0 & \cdots & \cdots & \theta_{s}%
\end{array}
\right)
\end{equation}
allowing us to express the vector $e^{i\mathbf{\theta}}$ in matrix form as%
\begin{equation}
e^{i\mathbf{\theta}}=e^{i\mathbf{\hat{\theta}}}\left(
\begin{array}
[c]{c}%
1\\
\vdots\\
\vdots\\
1
\end{array}
\right)  =\left(
\begin{array}
[c]{cccc}%
1 & 0 & \cdots & 0\\
\vdots & e^{i\theta_{2}} &  & \vdots\\
\vdots &  & \ddots & 0\\
0 & \cdots & \cdots & e^{i\theta_{s}}%
\end{array}
\right)  \left(
\begin{array}
[c]{c}%
1\\
\vdots\\
\vdots\\
1
\end{array}
\right)  .
\end{equation}
This gives
\begin{equation}
f\left(  \mathbf{\theta}\right)  =\left(
\begin{array}
[c]{ccc}%
1 & \cdots & 1
\end{array}
\right)  e^{-i\mathbf{\hat{\theta}}}\mathbb{M}e^{i\mathbf{\hat{\theta}}%
}\left(
\begin{array}
[c]{c}%
1\\
\vdots\\
1
\end{array}
\right)
\end{equation}
allowing us to obtain the Taylors series expansion about the point
$\mathbf{\theta}_{0}$ as
\begin{align}
f\left(  \mathbf{\theta}_{0}+\delta\mathbf{\theta}\right)   &  =f\left(
\mathbf{\theta}_{0}\right)  +\partial_{\mathbf{\theta}}f\left(  \mathbf{\theta
}_{0}\right)  ^{t}\delta\mathbf{\theta}+\delta\mathbf{\theta}^{t}%
\partial_{\mathbf{\theta}}^{2}f\left(  \mathbf{\theta}_{0}\right)
\delta\mathbf{\theta+\cdots}\nonumber\\
&  =\left(
\begin{array}
[c]{ccc}%
1 & \cdots & 1
\end{array}
\right)  e^{-i\widehat{\left(  \mathbf{\theta}_{0}+\delta\mathbf{\theta
}\right)  }}\mathbb{M}e^{i\widehat{\left(  \mathbf{\theta}_{0}+\delta
\mathbf{\theta}\right)  }}\left(
\begin{array}
[c]{c}%
1\\
\vdots\\
1
\end{array}
\right)  .\label{Taytheta}%
\end{align}
Noting that
\begin{equation}
e^{i\widehat{\left(  \mathbf{\theta}_{0}+\delta\mathbf{\theta}\right)  }%
}=e^{i\widehat{\mathbf{\theta}_{0}}}e^{i\widehat{\delta\mathbf{\theta}}%
}=e^{i\widehat{\mathbf{\theta}_{0}}}\left\{  \mathbb{I}_{s}+i\delta
\mathbf{\hat{\theta}}+i^{2}\left(  \delta\mathbf{\hat{\theta}}\right)
^{2}+\cdots\right\}  ,
\end{equation}
eq. $\left(  \ref{Taytheta}\right)  $gives
\begin{align}
f\left(  \mathbf{\theta}_{0}+\delta\mathbf{\theta}\right)   &  =\left(
\begin{array}
[c]{ccc}%
1 & \cdots & 1
\end{array}
\right)  \left\{  \mathbb{I}_{s}+\left(  -i\right)  \delta\mathbf{\hat{\theta
}}+\left(  -\frac{i}{2}\right)  ^{2}\left(  \delta\mathbf{\hat{\theta}%
}\right)  ^{2}+\cdots\right\}  \mathbb{M}\left(  \mathbf{\theta}\right)
\times\nonumber\\
&  \left\{  \mathbb{I}_{s}+i\delta\mathbf{\hat{\theta}}+\left(  \frac{i}%
{2}\right)  ^{2}\left(  \delta\mathbf{\hat{\theta}}\right)  ^{2}%
+\cdots\right\}  \left(
\begin{array}
[c]{c}%
1\\
\vdots\\
1
\end{array}
\right)  ,
\end{align}
where
\begin{equation}
\mathbb{M}\left(  \mathbf{\theta}\right)  =e^{-i\widehat{\mathbf{\theta}_{0}}%
}\mathbb{M}e^{i\widehat{\mathbf{\theta}_{0}}},
\end{equation}
showing that
\begin{align}
&  \delta^{1}f\left(  \mathbf{\theta}_{0}\right)  =i\left\{  \left(
\begin{array}
[c]{ccc}%
1 & \cdots & 1
\end{array}
\right)  \mathbb{M}\left(  \mathbf{\theta}\right)  \delta\mathbf{\hat{\theta}%
}\left(
\begin{array}
[c]{c}%
1\\
\vdots\\
1
\end{array}
\right)  -\left(
\begin{array}
[c]{ccc}%
1 & \cdots & 1
\end{array}
\right)  \delta\mathbf{\hat{\theta}}\mathbb{M}\left(  \mathbf{\theta}\right)
\left(
\begin{array}
[c]{c}%
1\\
\vdots\\
1
\end{array}
\right)  \right\}  \nonumber\\
&  =i\left\{  \left(
\begin{array}
[c]{ccc}%
1 & \cdots & 1
\end{array}
\right)  \mathbb{M}\left(  \mathbf{\theta}\right)  \delta\mathbf{\hat{\theta}%
}\left(
\begin{array}
[c]{c}%
1\\
\vdots\\
1
\end{array}
\right)  -\left[  \left(
\begin{array}
[c]{ccc}%
1 & \cdots & 1
\end{array}
\right)  \mathbb{M}\left(  \mathbf{\theta}\right)  \delta\mathbf{\hat{\theta}%
}\left(
\begin{array}
[c]{c}%
1\\
\vdots\\
1
\end{array}
\right)  \right]  ^{\ast}\right\}  \nonumber\\
&  =-2\left(
\begin{array}
[c]{ccc}%
1 & \cdots & 1
\end{array}
\right)  \operatorname{Im}\left\{  \mathbb{M}\left(  \mathbf{\theta}\right)
\right\}  \left(
\begin{array}
[c]{c}%
\delta\theta_{1}\\
\vdots\\
\delta\theta_{s}%
\end{array}
\right)
\end{align}
Hence
\begin{equation}
\partial_{\mathbf{\theta}}f\left(  \mathbf{\theta}_{0}\right)  ^{t}=-2\left(
\begin{array}
[c]{ccc}%
1 & \cdots & 1
\end{array}
\right)  \operatorname{Im}\left\{  \mathbb{M}\left(  \mathbf{\theta}\right)
\right\}  .
\end{equation}
The second order variation $\delta^{2}f\left(  \mathbf{\theta}_{0}\right)  $
is given by
\begin{align}
\delta^{2}f\left(  \mathbf{\theta}_{0}\right)   &  =-\frac{1}{2}\left(
\begin{array}
[c]{ccc}%
1 & \cdots & 1
\end{array}
\right)  \left\{  \overset{}{\mathbb{M}\left(  \mathbf{\theta}\right)  \left(
\delta\mathbf{\hat{\theta}}\right)  ^{2}+\left(  \delta\mathbf{\hat{\theta}%
}\right)  ^{2}\mathbb{M}\left(  \mathbf{\theta}\right)  }+\delta
\mathbf{\hat{\theta}}\mathbb{M}\left(  \mathbf{\theta}\right)  \delta
\mathbf{\hat{\theta}}\right\}  \left(
\begin{array}
[c]{c}%
1\\
\vdots\\
1
\end{array}
\right)  \nonumber\\
&  =\frac{1}{2}\left(
\begin{array}
[c]{ccc}%
1 & \cdots & 1
\end{array}
\right)  \left[  \delta\mathbf{\hat{\theta},}\left[  \mathbb{M}\left(
\mathbf{\theta}\right)  ,\delta\mathbf{\hat{\theta}}\right]  \right]  \left(
\begin{array}
[c]{c}%
1\\
\vdots\\
1
\end{array}
\right)
\end{align}
which gives
\begin{align}
& \left(
\begin{array}
[c]{ccc}%
1 & \cdots & 1
\end{array}
\right)  \mathbb{M}\left(  \mathbf{\theta}\right)  \left(  \delta
\mathbf{\hat{\theta}}\right)  ^{2}\left(
\begin{array}
[c]{c}%
1\\
\vdots\\
1
\end{array}
\right)  =\sum_{jk}M_{jk}\delta\theta_{k}^{2}\\
& \Rightarrow\frac{\partial}{\partial\theta_{l}}=2\sum_{j}M_{jl}\delta
\theta_{l}\\
& \left(
\begin{array}
[c]{ccc}%
1 & \cdots & 1
\end{array}
\right)  \left(  \delta\mathbf{\hat{\theta}}\right)  ^{2}\mathbb{M}\left(
\mathbf{\theta}\right)  \left(
\begin{array}
[c]{c}%
1\\
\vdots\\
1
\end{array}
\right)  =2\sum_{j}M_{lj}\delta\theta_{l}\\
& \left(
\begin{array}
[c]{ccc}%
1 & \cdots & 1
\end{array}
\right)  \delta\mathbf{\hat{\theta}}\mathbb{M}\left(  \mathbf{\theta}\right)
\delta\mathbf{\hat{\theta}}\left(
\begin{array}
[c]{c}%
1\\
\vdots\\
1
\end{array}
\right)  =\sum_{j}\left\{  M_{lj}\delta\theta_{j}+M_{jl}\delta\theta
_{j}\right\}
\end{align}%
\begin{equation}
f\left(  \mathbf{\theta}\right)  =\sum_{1\leq j,k\leq s}e^{-i\theta_{k}}%
M_{kj}e^{i\theta_{j}},
\end{equation}%
\begin{align}
\frac{\partial f\left(  \mathbf{\theta}\right)  }{\partial\theta_{l}}  &
=i\sum_{1\leq j\leq s}\left\{  e^{-i\theta_{j}}M_{jl}e^{i\theta_{l}%
}-e^{-i\theta_{l}}M_{lj}e^{i\theta_{j}}\right\}  \\
\frac{\partial^{2}f\left(  \mathbf{\theta}\right)  }{\partial\theta
_{l}\partial\theta_{m}}  & =-\left\{  \sum_{1\leq j\leq s}\left\{
e^{-i\theta_{j}}M_{jl}e^{i\theta_{l}}+e^{-i\theta_{l}}M_{lj}e^{i\theta_{j}%
}\right\}  \delta_{lm}-e^{-i\theta_{m}}M_{ml}e^{i\theta_{l}}-e^{-i\theta_{l}%
}M_{lm}e^{i\theta_{m}}\right\}
\end{align}%
\begin{equation}
\delta\mathbf{\theta=-}\left[  \frac{\partial^{2}f\left(  \mathbf{\theta
}\right)  }{\partial\mathbf{\theta}\partial\mathbf{\theta}}\right]  ^{-1}%
\frac{\partial f\left(  \mathbf{\theta}\right)  }{\partial\mathbf{\theta}}%
\end{equation}



\end{document}